% This table is the paper's primary benchmark reference.
% Overlap weights (Li et al. 2018) and Crump trimming (Crump et al. 2009) are
% the two external baselines against which Stabilized RIW (CS-DSOS, ours) is
% compared throughout the experiments.
\begin{table}[t]
  \centering
  \small
  \caption{Weighting approaches for common-support testing. Every method addresses overlap, but they
    differ in \emph{which quantity} they target and \emph{what} they preserve.
    The overlap-vs-stabilized contrast is central: both produce continuous, bounded weights,
    but the mechanism and estimand diverge fundamentally.}
  \label{tab:weighting-benchmarks}
  \begin{tabular}{lp{4.8cm}p{4.8cm}p{4.8cm}}
    \toprule
    & \textbf{Overlap weights}
    & \textbf{Crump trimming}
    & \textbf{Stabilized RIW (ours)} \\
    & (Li et al. 2018, 2019)
    & (Crump et al. 2009)
    & (Yamada et al. 2013) \\
    \midrule
    \multicolumn{4}{l}{\textit{What drives the weight}} \\
    \hspace{1em} Operates on
    & Propensity score $p(x)$
    & Propensity score $p(x)$
    & Density ratio $r(x) = \frac{p}{1-p} \cdot \frac{n_s}{n_t}$ \\
    \hspace{1em} Formula
    & $w(x) = p(x)\bigl(1-p(x)\bigr)$
    & $\mathbb{1}\{\min(p,1-p) \ge \alpha\}$
    & $w_{\mathrm{src}}\!=\!\frac{r}{(1-\lambda)+\lambda r}$,\;
      $w_{\mathrm{tgt}}\!=\!\frac{1}{\lambda+(1-\lambda)r}$ \\
    \hspace{1em} Prior ratio
    & Not reflected
    & Not reflected
    & Inferred from group sizes \\
    \midrule
    \multicolumn{4}{l}{\textit{Weight behavior}} \\
    \hspace{1em} Weight at $p\to0$
    & $\to 0$
    & Discarded if $p < \alpha$
    & $\to 0$ (source) \; $\to 1/\lambda$ (target) \\
    \hspace{1em} Weight at $p\to1$
    & $\to 0$
    & Discarded if $p > 1-\alpha$
    & $\to 1/\lambda$ (source) \; $\to 0$ (target) \\
    \hspace{1em} Bounded?
    & Yes ($\max = 0.25$)
    & N/A (binary)
    & Yes ($\max = 1/\lambda$) \\
    \hspace{1em} Zero at $p\in\{0,1\}$?
    & Yes (soft zero)
    & Yes (hard discard)
    & No (asymptotic, never zero) \\
    \midrule
    \multicolumn{4}{l}{\textit{Estimand and workflow}} \\
    \hspace{1em} Target population
    & Overlap region ($p(x)\approx 0.5$)
    & $p(x)\in[\alpha,\,1-\alpha]$
    & Full populations restricted to common support \\
    \hspace{1em} Tuning parameter
    & None (fixed)
    & $\alpha$ (arbitrary threshold)
    & $\lambda$ (principled via ESS) \\
    \hspace{1em} Preserves all observations?
    & Yes (soft downweight)
    & No (hard discard)
    & Yes (soft downweight) \\
    \hspace{1em} Under two-sided mismatch
    & Over-rejects
    & Over-rejects
    & Calibrated \\
    \bottomrule
  \end{tabular}

  \vspace{0.5em}
  \small
  \textbf{Why overlap and stabilized RIW look similar but differ.}
  Both produce \emph{continuous, bounded} weights that downweight extreme $p(x)$ values,
  and a casual reader may conflate them as two soft alternatives to Crump's hard trim.
  The separation is in what drives the weight:
  overlap weights target $p(1-p)$ --- the variance of a Bernoulli group indicator --- which
  peaks at $p=0.5$ and vanishes at the boundaries.
  Stabilized RIW targets the \emph{density ratio} $r(x)$, which encodes relative density of
  the two groups via Bayes' theorem, then bounds it through the sigmoid-like map
  $r \mapsto r/((1-\lambda)+\lambda r)$.
  Two consequences follow:
  (i)~overlap weights fix the weight shape by $p$ alone, whereas RIW incorporates
  the prior ratio and group-conditional densities --- an observation with $p=0.9$ gets
  overlap weight $0.09$ regardless of group sizes, but its stabilized RIW weight depends
  on whether it comes from the sparse or dense group;
  (ii)~overlap weights assign \emph{zero} weight at $p\in\{0,1\}$, a soft version of
  Crump's hard discard, while stabilized RIW preserves non-zero weight for every
  observation (bounded by $1/\lambda$), never completely discarding any data point.
  The empirical consequence is that overlap weights over-reject under two-sided weak-overlap
  contamination (Section~\ref{sec:calibration}), while the doubly weighted test remains
  calibrated because its asymmetric forward/inverse corrections handle each group's
  density structure separately.
\end{table}
